%==============================================================
%  Apresentação Beamer — Modelagem estatística de extremos espaciais
%  Autor: Enzo Porto Brasil   •   UnB
%==============================================================

\documentclass[aspectratio=169]{beamer}

% referências 
\bibliographystyle{apalike} % ou outra de sua preferência
\nocite{*}

% EXTRAS ------------------
\usepackage{comment} % para usar \begin{comment}
\usepackage{changepage}  % modificar os blocks
\usepackage{pdflscape}   % ou lscape
\usepackage{booktabs}    % para \toprule etc.
\usepackage{array}       % p/ colunas p{⟨largura⟩}
\usepackage{xcolor}      % p/ colorir símbolo concluído
\usepackage{pifont}      % p/ \ding{110} (■)
\usepackage{graphicx}
\usepackage{caption}      % \captionof

\usepackage{multirow}
\usepackage{float}      % permite usar [H] (fixa a tabela na posição exata)
\usepackage{amsmath}    % para usar expressões matemáticas em modo texto ($...$, \times, ^{-3}, etc.)
\usepackage{siunitx}    % opcional, mas ótimo para unidades físicas e alinhamento numérico
\usepackage{geometry}   % opcional, para ajustar margens caso precise de espaço extra

%--------------------------

\usepackage[brazil]{babel}
\addto\captionsbrazil{
  \renewcommand{\figurename}{Figura}
  \renewcommand{\tablename}{Tabela}
  \renewcommand{\partname}{Parte}
  \renewcommand{\chaptername}{Capítulo}
  \renewcommand{\appendixname}{Apêndice}
}
\setbeamertemplate{caption}[numbered]

% ---------- TEMA E CORES ----------
\usetheme{Madrid}
\definecolor{unbblue}{RGB}{0,38,100}
\definecolor{lightgray}{gray}{0.94}
\setbeamercolor{structure}{fg=unbblue}
\setbeamercolor{palette primary}{bg=unbblue, fg=white}
\setbeamercolor{palette secondary}{bg=unbblue, fg=white}
\setbeamercolor{palette tertiary}{bg=unbblue, fg=white}
\setbeamercolor{palette quaternary}{bg=unbblue, fg=white}
%\setbeamercolor{background canvas}{bg=lightgray} % basta alterar aqui se quiser fundo cinza claro

% ---------- METADADOS ----------
\title[Joint extremes of solar irradiance and air density]{%
  Joint extremes of solar irradiance and air density: \\ an EVT and copula approach to tail uncertainty}

\author[BRASIL, E. P]{\large{Author: Enzo Porto Brasil \\[-0.35cm] }}

\institute[]{\\[1.2cm] 
             \small{Department of Statistics, University of Brasília, Brazil} \\[-0.45cm]}

\date[]{\small{sys2025 - Huntsville, Alabama - November 15th, 2025}}


% ---------- PACOTES NECESSÁRIOS ----------
\usepackage{tikz}
\usetikzlibrary{calc}

% ---------- BOLINHAS DE PROGRESSO ----------
\newcommand{\progresscircles}{%
  \begin{tikzpicture}[remember picture,overlay]
    \foreach \i in {1,...,4} {
      \path let \p1 = ($(current page.north west)+(\i*0.5cm,-1cm)$) in
        node[circle,fill=gray!40,inner sep=1.5pt] at (\p1) {};
    }
    \path let \p2 = ($(current page.north west)+(\insertsectionnumber*0.5cm,-1cm)$) in
      node[circle,fill=unbblue,inner sep=1.5pt] at (\p2) {};
  \end{tikzpicture}%
}

\addtobeamertemplate{frametitle}{}{
  \vspace{-1.2em}
  \progresscircles
}

% ---------- CONTROLE DO RODAPÉ: MOSTRAR OU NÃO SEÇÃO ----------
\newif\ifshowsectionfooter
\showsectionfootertrue % padrão: mostrar seção

% ---------- PERSONALIZAÇÃO DA BARRA INFERIOR ----------
\setbeamertemplate{footline}{%
  \leavevmode
  \hbox{%
    \begin{beamercolorbox}[wd=\paperwidth,ht=2.5ex,dp=1ex,
                           leftskip=1em,rightskip=1em]{palette primary}
      \usebeamerfont{title in head/foot}%
      % ESQUERDA: SEÇÃO : SUBSEÇÃO
      \ifshowsectionfooter
        \ifx\insertsectionhead\empty
        \else
          \uppercase{\insertsectionhead}%
          \ifx\insertsubsectionhead\empty
          \else
            \,/\,\insertsubsectionhead
          \fi
        \fi
      \fi
      % DIREITA: título curto + nº de slide
      \hfill
      \insertshorttitle \quad \insertframenumber/\inserttotalframenumber
    \end{beamercolorbox}%
  }%
  \vskip0pt% <- cola a barra no rodapé real
}

% informa ao Beamer a altura ocupada pelo rodapé
% \setbeamersize{footer height=2.5ex}


% ---------- MARGENS INTERNAS ----------
\setbeamersize{text margin left=1.5cm, text margin right=1.5cm}

% ---------- SUMÁRIO AO INICIAR SEÇÃO ----------
\AtBeginSection[]{
  \begin{frame}{Agenda}
    \tableofcontents[currentsection]
    \vfill
  \end{frame}
}


% INÍCIO DO DOCUMENTO
\begin{document}

% CAPA --------------------------------------------------------
% Substitua \maketitle por:
% Orientadora: Prof. Cira E. G. Otiniano
\begin{frame}
  \titlepage
\end{frame}



% SUMÁRIO GERAL ----------------------------------------------
{\begin{frame}{Agenda}
  \tableofcontents
\end{frame}
}

\section{Introduction}

%%%%%%%%%%%%%%%%%%%%%%%%%%%%%%%%%%%%%%%%

\begin{comment}
\begin{frame}{News}
  \centering

  %  Ajuste: tira 1 cm de cada lado ─ mexa nesses números até “encher” de fato
  \begin{adjustwidth*}{-1.5cm}{-1.5cm}

  % Bloco horizontal com largura exata da lâmina
  \makebox[\paperwidth][c]{%
    \includegraphics[width=0.32\paperwidth,
                     page=1,trim=0 5 0 15,clip]{Noticias extremas/noticia_1.pdf}%
    \hspace{0.008\paperwidth}%
    \includegraphics[width=0.32\paperwidth,
                     page=1,trim=0 5 0 15,clip]{Noticias extremas/noticia_2.pdf}%
    \hspace{0.008\paperwidth}%
    \includegraphics[width=0.32\paperwidth,
                     page=1,trim=0 5 0 15,clip]{Noticias extremas/noticia_3.pdf}%
  }

  \vspace{-0.3cm}  % pequeno espaço entre imagem e legendas

  % Legendas alinhadas horizontalmente sob cada imagem
  {\scriptsize
    \hspace*{0.1\paperwidth} Al Jazeera, 2025
    \hspace*{0.15\paperwidth} European Commission, 2025
    \hspace*{0.13\paperwidth} BBC News, 2025
  }
\end{adjustwidth*}
\end{frame}
\end{comment}

\begin{frame}{Context}
  \centering
  \includegraphics[
    width=\textwidth,
    height=0.85\textheight,
    keepaspectratio
  ]{Noticias extremas/news1.pdf}
\end{frame}

\begin{frame}{Context}
  \centering
  \includegraphics[
    width=\textwidth,
    height=0.85\textheight,
    keepaspectratio
  ]{Noticias extremas/news2.pdf}
\end{frame}

\begin{frame}{Extreme events}
  \centering
  \includegraphics[
    width=\textwidth,
    height=0.85\textheight,
    keepaspectratio
  ]{Noticias extremas/extremes.png}
\end{frame}


%%%%%%%%%%%%%%%%%%%%%%%%%%%%%%%%%%%%%%%%

\begin{comment}
%%%%%%%%%%%%%%%%%%%%%%%%%%%%%%%%%%%%%%%%
\begin{frame}{Motivation}
  \centering

  %  Ajuste: tira 1 cm de cada lado ─ mexa nesses números até “encher” de fato
  \begin{adjustwidth*}{-1.5cm}{-1.5cm}

  % Bloco horizontal com largura exata da lâmina
  \makebox[\paperwidth][c]{%
    \includegraphics[width=0.32\paperwidth,
                     page=1,trim=0 5 0 15,clip]{Noticias extremas/noticia_4.pdf}%
    \hspace{0.008\paperwidth}%
    \includegraphics[width=0.32\paperwidth,
                     page=1,trim=0 5 0 15,clip]{Noticias extremas/noticia_5.pdf}%
    \hspace{0.008\paperwidth}%
    \includegraphics[width=0.32\paperwidth,
                     page=1,trim=0 5 0 15,clip]{Noticias extremas/noticia_6.pdf}%
  }

  \vspace{-0.3cm}  % pequeno espaço entre imagem e legendas

  % Legendas alinhadas horizontalmente sob cada imagem
  {\scriptsize
    \hspace*{0.045\paperwidth} NASA Earth Observatory, 2025
    \hspace*{0.09\paperwidth} The New York Times, 2025
    \hspace*{0.14\paperwidth} SciTechDaily, 2025
  }
\end{adjustwidth*}
\end{frame}



%%%%%%%%%%%%%%%%%%%%%%%%%%%%%%%%%%%%%%%%
\begin{frame}{Historical context}

\begin{enumerate}
    \item Fisher and Tippett (1928): pioneering results in Extreme Value Theory (EVT)
        \vspace*{0.7cm}
    \item Von Mises (1936): fundamental contributions to the theoretical foundations of EVT
        \vspace*{0.7cm}        
    \item Gnedenko (1943): rigorous formalization of the limiting extreme value distributions
        \vspace*{0.7cm}
    \item Jenkinson (1955): establishment of the Generalized Extreme Value (GEV) distribution
        \vspace*{0.7cm}   
    \item 1940s–1955 to present: GEV as the primary choice for modeling univariate extremes
        \vspace*{0.7cm}   
    \item Recognized limitations of the GEV under complex or multimodal scenarios
\end{enumerate}


% refs
\begin{flushright}
{\tiny Referências: Haan e Ferreira, 2006; Davison, Padoan e Ribatet, 2012; Kotz e Nadarajah, 2000; Otiniano \textit{et al.}, 2023b}
\end{flushright}


\end{frame}


%%%%%%%%%%%%%%%%%%%%%%%%%%%%%%%%%%%%%%%%
\begin{frame}{Contexto histórico}
  \setbeamercovered{transparent}

  \begin{enumerate}
  \setcounter{enumi}{5}
    \item Surgem várias extensões e generalizações da GEV, por exemplo:\\
    \begin{itemize}
      \item Aryal e Tsokos (2009): Transmutada GEV (TGEV)
      \item Provost \textit{et al.} (2018): q-GEV
      \item Gramosa \textit{et al.} (2020): GEV inflada de zeros (ZIGEV)
      \item Otiniano \textit{et al.} (2023): GEV Bimodal (BGEV)
    \end{itemize}
  \end{enumerate}

  \pause
  \vspace*{0.5cm}

  A \textbf{BGEV} de Otiniano \textit{et al.} (2023):
  \begin{itemize}
    \item Tem um parâmetro a mais;
    \item Apresenta mais flexibilidade que a GEV; e
    \item Captura dados extremos com heterogeneidade e múltiplos regimes
  \end{itemize}

\end{frame}
\end{comment}




\section{Theoretical background}

%**********************************************
\subsection{Extreme Value Theory (EVT)}
%**********************************************
\begin{frame}{Extreme Value Theory (EVT)}

\small
\textbf{The interest goes beyond the general behavior of the data:} \\
the focus now is directly on \textit{tail uncertainty.}


\vspace{1cm}

\textbf{How can we model the heavy tails of the data?}

\begin{itemize}
  \item Stochastic behavior of minima and maxima in i.i.d. sequences
  \item Blocks maxima \& minima: partition data into blocks to extract extremes
  \item GEV and Bimodal GEV distributions: asymptotic limits as $n \to \infty$
  \item Risk assessment: Value at Risk (VaR), return levels, and extreme forecasts
\end{itemize}


\vfill
\begin{flushright}
{\tiny References: Fisher and Tippett (1928);  Von Mises (1936); Gnedenko (1943); Kotz and Nadarajah (2000); Yuliang et al. (2021)}
\end{flushright}
\end{frame}




%%%%%%%%%%%%%%%%%%%%%%%%%%%%%%%%%%%%%%%%
\begin{frame}{Generalized Extreme Value distribution (GEV)}

$\boldsymbol{Y \sim \mathbf{GEV}(\cdot;\,\xi,\mu,\sigma)}$

\vspace*{0.3cm}
\begin{block}{GEV density}
  \resizebox{\linewidth}{!}{$
    f_{\text{GEV}}(y;\mu,\sigma,\xi)=
    \begin{cases}
      \dfrac{1}{\sigma}\bigl[1+\xi \tfrac{(y-\mu)}{\sigma}\bigr]^{-1/\xi-1}
      \exp\!\Bigl\{-\bigl[1+\xi \tfrac{(y-\mu)}{\sigma}\bigr]^{-1/\xi}\Bigr\},
      & \xi\neq0,\\[14pt]
      \dfrac{1}{\sigma}\,
      \exp\!\Bigl\{-\tfrac{(y-\mu)}{\sigma}-\exp[-\tfrac{(y-\mu)}{\sigma}]\Bigr\},
      & \xi=0.
    \end{cases}
  $}
\end{block}

\vspace*{0.5cm}
\noindent Where $\xi \in \mathbb{R}$ (shape), $\mu \in \mathbb{R}$ (location), and $\sigma > 0$ (scale). With support given by 
$1+\xi\,(y-\mu)/\sigma>0$ in general, i.e., $y>\mu-\sigma/\xi$ if $\xi>0$ and $y<\mu-\sigma/\xi$ if $\xi<0$, and $y\in\mathbb{R}$ when $\xi=0$.

\vfill
\begin{flushright}
{\tiny References: Von Mises(1936); Jenkinson (1955); Kotz and Nadarajah (2000); e Longin (2016)}
\end{flushright}
\end{frame}


\begin{comment}
\vspace*{0.25cm}
A distribuição GEV assume diferentes formas dependendo do valor de \(\xi\), e se reduz às distribuições:\\
\begin{center} Gumbel (\(\xi = 0\)); Fréchet (\(\xi > 0\)); e Weibull negativa (\(\xi < 0\)).\end{center}
\end{comment}


%%%%%%%%%%%%%%%%%%%%%%%%%%%%%%%%%%%%%%%%%
\begin{frame}{Generalized Extreme Value distribution (GEV)}
\begin{figure}[ht]
    \centering
    \includegraphics[width=\textwidth]{distribuicao GEV.pdf}
    \caption*{{\usebeamercolor[fg]{caption name} \bfseries Figure.} Probability density function and cumulative distribution function of the GEV with $\mu=0$, $\sigma=2$, and $\xi$ varying as indicated in the legend}
\end{figure}



\end{frame}

%%%%%%%%%%%%%%%%%%%%%%%%%%%%%%%%%%%%%%%%%
\begin{frame}{Bimodal Generalized Extreme Value distribution (BGEV)}

\vspace*{0.45cm}
$\boldsymbol{X \sim \mathbf{BGEV}(\cdot;\xi,\mu,\sigma,\delta)}$

\vspace*{0.3cm}
\begin{block}{BGEV density}
    \resizebox{\linewidth}{!}{$
    \displaystyle
    f_{\text{BGEV}}(x;\xi,\mu,\sigma,\delta)=
    \begin{cases}
    \dfrac{1}{\sigma}\Bigl[1+\xi \tfrac{T(x)}{\sigma}\Bigr]^{-1/\xi-1}
            \exp\!\Bigl\{-\Bigl[1+\xi \tfrac{T(x)}{\sigma}\Bigr]^{-1/\xi}\Bigr\}
            \,T'(x), & \xi\neq0,\\[14pt]
    \dfrac{1}{\sigma}\,
            \exp\!\bigl(-\tfrac{T(x)}{\sigma}\bigr)\;
            \exp\!\Bigl\{ -\exp\!\bigl(-\tfrac{T(x)}{\sigma}\bigr)\Bigr\}
            \,T'(x), & \xi=0.
    \end{cases}
    $}
\end{block}

\vspace*{0.5cm}

Where $\xi \in \mathbb{R}$, $\delta > -1$, $\sigma > 0$ (shape parameters), and $\mu \in \mathbb{R}$ (location parameter). Additionally, $T(x) = (x - \mu)\, |x - \mu|^{\delta}$ and $T'(x) = (\delta + 1)\, |x - \mu|^{\delta}$.

\vspace*{0.25cm}
\begin{itemize}
\item Note: when $\delta = 0$, $\boldsymbol{X \sim \text{\textbf{BGEV}}(\, \xi, \, \mu, \, \sigma, \, 0) = \text{\textbf{GEV}}(\xi, \mu, \sigma)}$
\end{itemize}

\vfill
\begin{flushright}
{\tiny References: Otiniano et al. (2023 and 2025)}
\end{flushright}

\end{frame}


\begin{comment} % SUPORTE DA BGEV
    \begin{equation*}
    \label{bgev4_suporte_densidade}
    \mathrm{supp}\left( f_{\text{BGEV}}(\cdot; \xi, \mu, \sigma, \delta) \right) =
    \begin{cases}
    \left[ \mu - \left( \frac{\sigma}{\xi} \right)^{\frac{1}{\delta + 1}}, +\infty \right), & \xi > 0, \\
    \left( -\infty, \mu + \left| \frac{\sigma}{\xi} \right|^{\frac{1}{\delta + 1}} \right], & \xi < 0, \\
    (-\infty, +\infty), & \xi = 0.
    \end{cases}
    \end{equation*}


%%%%%%%%%%%%%%%%%%%%%%%%%%%%%%%%%%%%%%%%%
\begin{frame}{Distribuição Generalizada de Valor Extremo Bimodal (BGEV)}

\begin{itemize}
\item Se $\delta = 0$: $\boldsymbol{X \sim \text{\textbf{BGEV}}(\, \xi, \, \mu, \, \sigma, \, 0) = \text{\textbf{GEV}}(\xi, \mu, \sigma)}$

\vspace*{1cm}
\item  \textbf{Função quantílica}
\vspace*{-0.25cm}
    \begin{equation*}
    \label{bgev5_quantil}
    Q(x) =
    \begin{cases}
    \operatorname{sgn} \left\{ \dfrac{\sigma}{\xi} \left[ (-\log(x))^{-\xi} - 1 \right] \right\}
    \left| \dfrac{\sigma}{\xi} \left[ (-\log(x))^{-\xi} - 1 \right] \right|^{\frac{1}{\delta + 1}} + \mu, & \text{se } \xi \neq 0, \\[10pt]
    \operatorname{sgn} \left[ -\sigma \log(-\log(x)) \right] \left| -\sigma \log(-\log(x)) \right|^{\frac{1}{\delta + 1}} + \mu, & \text{se } \xi = 0.
    \end{cases}
    \end{equation*}
\end{itemize}

\vfill
\begin{flushright}
{\tiny References: Otiniano et al., 2023}
\end{flushright}


\end{frame}
\end{comment}



%%%%%%%%%%%%%%%%%%%%%%%%%%%%%%%%%%%%%%%%%%%%%%%
\begin{frame}{Bimodal Generalized Extreme Value distribution (BGEV)}

\begin{figure}[ht]
    \centering
    \includegraphics[width=\textwidth]{distribuicao BGEV.pdf}
    \caption*{{\usebeamercolor[fg]{caption name} \bfseries Figure.} 
    Probability density function and cumulative distribution function, respectively, 
    of the BGEV with $\mu = 0$, $\sigma = 1$, $\xi = 0$, and $\delta$ varying as indicated in the legend}
\end{figure}

\end{frame}


%**********************************************
\subsection{Copulas}
%**********************************************

\begin{comment}
\begin{frame}{}
\footnotesize
\textbf{Idea.} A copula $C:[0,1]^d \!\to [0,1]$ is a multivariate distribution function with U[0,1] margins, 
which captures the \emph{dependence structure} independently of the marginal behavior.

\textbf{Why use it?}
\begin{itemize}
  \item Build flexible multivariate models (choose margins + choose dependence)
  \item Capture tail (a)symmetry and tail dependence for risk/EVT
  \item Simulate joint VaR, return levels, stress scenarios
\end{itemize}

\vspace*{0.5cm}
\textbf{Sklar’s theorem.}  
For a joint CDF $H(x_1,\dots,x_d)$ with marginal distributions $F_1,\dots,F_d$, there exists a copula $C:[0,1]^d\!\to[0,1]$ such that
\[
H(x_1,\dots,x_d)
  = C\!\big(F_1(x_1),\dots,F_d(x_d)\big),
\]
and if the margins are continuous, $C$ is unique. Conversely, for any copula $C$ and margins $F_1,\dots,F_d$,
\[
C(u_1, u_2, \dots, u_d)
  = H\!\Bigl(F_1^{-1}(u_1), \dots, F_d^{-1}(u_d)\Bigr),
\]
where each $u_i = F_i(x_i) \in [0,1]$ represents the \textit{probability scale} value of variable $X_i$.

\vfill
\begin{flushright}
{\tiny References: Sklar (1959); Nelsen (2006); Trivedi and Zimmer (2007); Joe (2014)}
\end{flushright}

\end{frame}
\end{comment}

%%% %%% ALTERNATIVE %%% %%%
\begin{frame}{Copula Theory}
\footnotesize
\begin{columns}[T]
\begin{column}{0.48\textwidth}
\textbf{Idea.}  
A copula $C:[0,1]^d\!\to[0,1]$ is a multivariate distribution 
with Uniform(0,1) margins, capturing the dependence structure 
independently of the marginals.

\vspace{0.5cm}
\textbf{Why use it?}
\begin{itemize}
  \item Flexible modeling (margins + dependence)
  \item Tail (a)symmetry and dependence
  \item Joint risk and stress simulations
\end{itemize}
\end{column}

\begin{column}{0.52\textwidth}
\textbf{Sklar’s theorem.}
For a joint CDF $H(x_1,\dots,x_d)$ with marginal distributions $F_1,\dots,F_d$, there exists a copula $C:[0,1]^d\!\to[0,1]$ such that
\[
H(x_1,\dots,x_d)
  = C\!\big(F_1(x_1),\dots,F_d(x_d)\big),
\]
and if the margins are continuous, $C$ is unique. Conversely, for any copula $C$ and margins $F_1,\dots,F_d$,
\[
C(u_1, u_2, \dots, u_d)
  = H\!\Bigl(F_1^{-1}(u_1), \dots, F_d^{-1}(u_d)\Bigr),
\]
where each $u_i = F_i(x_i) \in [0,1]$ represents the \textit{probability scale} value of variable $X_i$.
\end{column}
\end{columns}

\vfill
\begin{flushright}
{\tiny References: Sklar (1959); Nelsen (2006); Trivedi and Zimmer (2007); Joe (2014)}
\end{flushright}

\end{frame}



%==================== SLIDE 2 ====================
\begin{frame}{Copulas: validity (bivariate case)}
\footnotesize

\begin{center}
\textbf{When is $C$ a bivariate copula?}
\end{center}

\vspace*{0.3cm}

Let $C:[0,1]^2 \to [0,1]$. Then $C$ is a copula if:

\begin{enumerate}
  \item $C$ is \textbf{non-decreasing} in each argument;
  \item $C$ is \textbf{right-continuous} in each argument;
  \item \textbf{Uniform margins on $[0,1]$:}
  \[
    C(u,1) = u,\quad C(1,v) = v,\quad
    C(u,0) = 0,\quad C(0,v) = 0,\quad C(1,1)=1;
  \]
  \item \textbf{Bivariate monotonicity}: for any rectangle 
  $[a_1,b_1]\times[a_2,b_2]\subset[0,1]^2$,
  \[
    C(b_1,b_2) - C(a_1,b_2) - C(b_1,a_2) + C(a_1,a_2) \ge 0.
  \]
\end{enumerate}

\vfill
\begin{flushright}
{\tiny References: Nelsen (2006); Trivedi and Zimmer (2007); Joe (2014)}
\end{flushright}
\end{frame}


\begin{frame}{Copulas: bivariate scenarios}
\begin{figure}[ht]
  \centering
  \setlength{\abovecaptionskip}{1.2pt}   % ↓ distância entre imagem e legenda
  \setlength{\belowcaptionskip}{0pt}   % ↓ distância após a legenda (opcional)
  \includegraphics[width=1\textwidth]{copulas_sim.pdf}
  \caption*{{\usebeamercolor[fg]{caption name}\bfseries Figure.} Examples of simulated bivariate copulas (Gaussian, Frank, Clayton, Joe) under two marginal settings: Normal-Normal and Lognormal-Gumbel}
\end{figure}


\end{frame}



%**********************************************
\subsection{Complex systems}
%**********************************************
\begin{frame}{Complex Systems: linking EVT and Copulas}
\small

\textbf{Idea.}  
Real phenomena show \textit{non-Gaussian extremes}: heavy tails, multimodality, nonlinear/assymmetric dependence, and nonstationarity.

\vspace{0.3cm}
\textbf{Modeling strategy.}
\begin{itemize}
  \item \textbf{Marginals:} GEV/BGEV for extremes;
  \item \textbf{Dependence:} Copulas
  \item \textbf{Nonstationary:} parameters as functions of covariates
\end{itemize}

\vspace{0.25cm}
\textbf{Workflow.}
\begin{itemize}
  \item Fit marginal tails $\Rightarrow$ transform to $u_i = F_i(y_i)$
  \item Fit copula on $(u_1,u_2)$ $\Rightarrow$ infer tail dependence, joint risk
\end{itemize}

\vspace{0.25cm}
\textbf{Key takeaway:}  
EVT captures how extreme the values are; Copulas capture how they are related. Together they describe the complexity of joint extremes.

\vfill
\begin{flushright}
{\tiny References: Kotz and Nadarajah (2000); Sornette (2004); Joe (2014)}
\end{flushright}

\end{frame}


\section{Material and methods}

\begin{frame}{Research objective}
Using Extreme Value Theory (EVT) and copula-based dependence models, this study aims to improve the robustness of tail inference beyond the traditional GEV framework, especially under complex data structures where standard assumptions fail. 
\vspace*{0.3cm}
In order to assess risk measures linking:
\begin{itemize}
\item[-] \textbf{Solar Spectral Irradiance (SSI)};
\item[-] \textbf{Absolute uncertainty of SSI measurements};
\item[-] \textbf{Air density}.
\end{itemize}
\end{frame}


\begin{frame}{Material}

\textbf{Datasets overview}
\vspace*{0.5cm}
\begin{itemize}
  \item \textbf{Solar Spectral Irradiance (SSI)} - NASA SORCE/XPS mission (photodiode 0.1–7 nm, 2005–2019), daily interpolated series [W/m$^2$];
    \vspace*{0.15cm}
  \item \textbf{Absolute Uncertainty (AU) of SSI measurements} - Instrumental 1$\sigma$ uncertainty [W/m$^2$] from the same mission and period for SSI;
    \vspace*{0.15cm}
  \item \textbf{Air Density (AD)} - Obtained from the Brazilian National Meteorological Institute (INMET) station in Brasília, using temperature, pressure, and humidity (Picard et al., 2008).
\end{itemize}
\end{frame}

\begin{frame}{Methods}

\textbf{Modeling steps}
\vspace*{0.5cm}
\begin{enumerate}
  \item \textbf{Univariate modeling}: GEV and BGEV distributions fitted to minima (364 blocks of 15 days) and maxima (90 blocks of 61 days);
  \vspace*{0.15cm}
  \item \textbf{Dependence structure}: parametric bivariate copulas fitted to extreme blocks of SSI, AI, and AD;
  \vspace*{0.15cm}
  \item \textbf{Risk measures}: ``bivariate'' Value at Risk (VaR) and return levels derived from the selected copulas.
\end{enumerate}


\end{frame}
\section{Main results}

%%%%%%%%%%%%%%%%%% %%%%%%%%%%%%%%%%%% 
\begin{frame}{Datasets}
\begin{figure}[ht]
    \centering
    % Slightly exceed text width for a broader visual span on the slide
    \makebox[\textwidth][c]{\includegraphics[width=1.12\textwidth]{dataset.png}}
    \caption*{{\usebeamercolor[fg]{caption name}\bfseries Figure.}\;
    Time series illustrating the Solar Spectral Irradiance from the Lyman-alpha model (left) and the datasets analyzed in this study (right): SORCE mission observations of SSI and Brasília’s daily air density extremes derived from INMET atmospheric variables}
\end{figure}

\end{frame}


%%%%%%%%%%%%%%%%%% %%%%%%%%%%%%%%%%%% 
\begin{frame}{Univariate modeling}
\begin{figure}[ht]
    \centering
    % Slightly exceed text width for a broader visual span on the slide
    \makebox[\textwidth][c]{\includegraphics[width=\textwidth]{univariate.pdf}}
    \caption*{{\usebeamercolor[fg]{caption name}\bfseries Figure.}\;
    Empirical density and univariate fitting models of the GEV and BGEV distributions applied to block of maxima and minima}
\end{figure}

\end{frame}


%%%%%%%%%%%%%%%%%% %%%%%%%%%%%%%%%%%% 
\begin{frame}{Dependence structure}
\begin{figure}[ht]
    \centering
    % Slightly exceed text width for a broader visual span on the slide
    \makebox[\textwidth][c]{\includegraphics[width=1\textwidth]{Copulas/correlacoes.png}}
    \caption*{{\usebeamercolor[fg]{caption name}\bfseries Figure.}\;
    Scatter plots with histograms and corresponding empirical densities for block maxima and minima, referring to the log-returns of Solar Spectral Irradiance (daily mean), air density, and absolute uncertainty}
\end{figure}
\end{frame}

\begin{frame}{Copula fitting: scenario 1}
\begin{center}\textbf{MAX: Log-return(Solar Spectral Irradiance) vs. Air Density}\end{center}

\begin{figure}[ht]
    \centering
    % Slightly exceed text width for a broader visual span on the slide
    \makebox[\textwidth][c]{\includegraphics[width=1.12\textwidth]{Copulas/ISxDA_max.pdf}}
    \caption*{{\usebeamercolor[fg]{caption name}\bfseries Figure.}\;
    Empirical joint density and three best-fitting copulas (Independent, Gaussian, and Frank - based on AIC) for \textbf{block maxima} of SSI log-returns (GEV model) and Air Density (BGEV model)}
\end{figure}
\end{frame}


\begin{frame}{Copula fitting: scenario 1}
\begin{center}\textbf{MAX: Log-return(Solar Spectral Irradiance) vs. Air Density}\end{center}
\begin{figure}[ht]
    \centering
    % Slightly exceed text width for a broader visual span on the slide
    \makebox[\textwidth][c]{\includegraphics[width=1.12\textwidth]{Copulas/3d - ISxDA_max.pdf}}
    \caption*{{\usebeamercolor[fg]{caption name}\bfseries Figure.}\;
    From left to right: fitted copula surface (Independent), empirical joint density surface, and their difference (model – empirical) for \textbf{block maxima} of SSI log-returns (GEV model) and Air Density (BGEV model)}
\end{figure}

\end{frame}

\begin{frame}{Copula fitting: scenario 2}
\begin{center}\textbf{MIN: Log-return(Solar Spectral Irradiance) vs. Air Density}\end{center}
\begin{figure}[ht]
    \centering
    % Slightly exceed text width for a broader visual span on the slide
    \makebox[\textwidth][c]{\includegraphics[width=1.12\textwidth]{Copulas/ISxDA_min.png}}
    \caption*{{\usebeamercolor[fg]{caption name}\bfseries Figure.}\;
    From left to right: empirical joint density with fitted copula (Tawn Type I, rotated $270^\circ$), copula surface, empirical joint density surface, and their difference (model – empirical) for \textbf{block minima} of SSI log-returns (GEV model) and Air Density (BGEV model)}
\end{figure}
\end{frame}



\begin{frame}{Copula fitting: scenario 3}
\begin{center}\textbf{MIN: Log-return(Solar Spectral Irradiance) vs. Absolute Uncertainty}\end{center}
\begin{figure}[ht]
    \centering
    % Slightly exceed text width for a broader visual span on the slide
    \makebox[\textwidth][c]{\includegraphics[width=1.12\textwidth]{Copulas/ISxAI_min.png}}
    \caption*{{\usebeamercolor[fg]{caption name}\bfseries Figure.}\;
    From left to right: empirical joint density with fitted copula (Frank), copula surface, empirical joint density surface, and their difference (model – empirical) for \textbf{block minima} of SSI log-returns (GEV model) and Absolute Uncertainty (BGEV model)}
\end{figure}

\end{frame}


\begin{comment}
\begin{frame}{Risk estimation: extremes of Solar Spectral Irradiance}

\begin{table}[ht]
\centering
\footnotesize
\caption{Expected daily values and percentage increase of the mean SSI (\(\times10^{-4}\)~W\,m\(^{-2}\)) for different return periods (\(T\)), considering minimum (15-day) and maximum (61-day) blocks. Overall mean irradiance: \(2.625\times10^{-4}\)~W\,m\(^{-2}\).}
\vspace{0.2cm}
\begin{tabular}{cc@{\hspace{0.3cm}}cc@{\hspace{0.6cm}}cc@{\hspace{0.3cm}}c}
\toprule
\multirow{2}{*}{$T_{\text{blocks}}$} &
\multirow{2}{*}{$T_{\text{days}}$} &
\multicolumn{2}{c}{\textbf{Minima}} &
\multicolumn{2}{c}{\textbf{Maxima}} \\
\cmidrule(lr){3-4} \cmidrule(lr){5-6}
& & $\text{SSI (daily)}$ & Increase (\%) & $\text{SSI (daily)}$ & Increase (\%) \\ 
\midrule
2   & 30   & 3.06 $\pm$ 0.04 & 16.8 & 3.67 $\pm$ 0.17 & 39.7 \\
5   & 75   & 3.56 $\pm$ 0.10 & 35.7 & 4.64 $\pm$ 0.38 & 76.8 \\
20  & 300  & 4.74 $\pm$ 0.41 & 80.8 & 6.79 $\pm$ 1.31 & 158.6 \\
50  & 750  & 6.13 $\pm$ 0.96 & 133.5 & 9.11 $\pm$ 2.83 & 246.9 \\
100 & 1500 & 7.82 $\pm$ 1.83 & 198.0 & 11.72 $\pm$ 4.90 & 346.4 \\
\bottomrule
\end{tabular}
\end{table}

\end{frame}
\end{comment}


\begin{frame}{Risk estimation: scenario 1}
\begin{columns}[T,onlytextwidth]
  \begin{column}{0.75\textwidth}
  \vspace*{-.37cm}
    \begin{adjustwidth*}{-1.3cm}{-1.3cm}  % aumenta área útil (ajuste o valor)
      \includegraphics[width=1.2\textwidth]{Copulas/quadro_geral_ISxDA_max.pdf}
    \end{adjustwidth*}
  \end{column}
  \begin{column}{0.25\textwidth}
    \vspace*{3.5cm}
    \scriptsize
    \textbf{Description.}\\[0.2em] Summary of the complete modeling workflow: from Solar Spectral Irradiance (SSI) and air density data to marginal and copula fitting, leading to the estimation of the Value at Risk (VaR) for \textbf{block maxima}.
  \end{column}
\end{columns}
\end{frame}



\begin{frame}{Risk estimation: scenario 2}
\begin{columns}[T,onlytextwidth]
  \begin{column}{0.75\textwidth}
  \vspace*{-.37cm}
    \begin{adjustwidth*}{-1.3cm}{-1.3cm}  % aumenta área útil (ajuste o valor)
      \includegraphics[width=1.2\textwidth]{Copulas/quadro_geral_ISxDA_min.pdf}
    \end{adjustwidth*}
  \end{column}
  \begin{column}{0.25\textwidth}
    \vspace*{3.5cm}
    \scriptsize
    \textbf{Description.}\\[0.2em] Summary of the complete modeling workflow: from Solar Spectral Irradiance (SSI) and air density data to marginal and copula fitting, leading to the estimation of the Value at Risk (VaR) for \textbf{block minima}.
  \end{column}
\end{columns}
\end{frame}

\section{Summary}


%*************************************
\begin{frame}{Summary}
\large
\begin{itemize}
\item EVT is essential to understand solar–atmospheric behavior.
\vspace*{0.55cm}
\item Extremes of solar irradiance remain highly uncertain, requiring EVT-based modeling.
\vspace*{0.55cm}
\item GEV/BGEV margins and copulas captured the main dependence structure.
\vspace*{0.55cm}
\item SSI, absolute uncertainty, and air density measured in Brasília show weak but detectable joint behavior.
\vspace*{0.55cm}
\item Systematic errors and data gaps must be treated with care in inference.
\end{itemize}
\end{frame}



%*************************************
\begin{frame}{Acknowledgments}
\begin{figure}[ht]
    \centering
    % Slightly exceed text width for a broader visual span on the slide
    \makebox[\textwidth][c]{\includegraphics[width=1\textwidth]{brasilia (1).png}}
    \caption*{\tiny Contributors: Fidel Morales, Enzo Brasil, Cira Otiniano, Carolyne Brito, Antonio Caio, Júlia Birbeire, Beatriz Albernaz, and Jéssica Ferreira}

\end{figure}

\end{frame}
\section*{References}

\begin{frame}[allowframebreaks]{Selected References}
\scriptsize
\bibliography{ref} % sem extensão
\end{frame}





%==============================================================
\end{document}
